\begin{abstract}
Although large unsupervised language models (LMs) acquire extensive world knowledge and some capability for reasoning, fine-grained manipulation of their outputs remains challenging due to the purely unsupervised training regime. Current techniques for imparting controllability generally involve assembling human judgments on the comparative quality of generated outputs, then refining the LM to adhere to these preferences—often utilizing reinforcement learning from human feedback (RLHF). However, RLHF is intricate and frequently suffers from instability, typically involving a two-stage process: first fitting a reward model to mirror human preferences, and then tuning the LM with reinforcement learning to optimize the predicted reward, all while constraining divergence from the original model. In this work, we propose a novel reparameterization of the reward model in the RLHF framework that facilitates closed-form retrieval of the corresponding optimal policy. This advancement reduces the RLHF procedure to straightforward classification, replacing the need for sampling during LM fine-tuning and extensive hyperparameter optimization. The proposed method, named \textit{Direct Preference Optimization} (DPO), is efficient, robust, and easy to use, greatly simplifying the training pipeline. Empirically, we demonstrate that DPO aligns LMs to human preferences on par with or surpassing the performance of prior approaches. Remarkably, DPO-based fine-tuning surpasses PPO-style RLHF for sentiment modulation in generated text and either matches or exceeds state-of-the-art outcomes in summarization and single-turn conversation, all while remaining considerably easier to apply.
\end{abstract}

\section{Introduction}
Large-scale unsupervised language models (LMs) trained on massive datasets have demonstrated a range of unexpected competencies~\citep{chowdhery2022palm, brown2020language, touvron2023llama,bubeck2023sparks}. Nonetheless, because the training data originates from humans whose objectives, motivations, and expertise vary widely, these models absorb a broad spectrum of behaviors. Some aspects of this human-derived data are not suitable for direct emulation; for instance, while it is advantageous for an AI coding assistant to \textit{recognize} common programming errors so it can provide corrections, we would prefer it to generate code that reflects the (often scarce) high-level expertise present in its training set, rather than simply replicating the frequency of common mistakes. Analogously, we desire that language models be \textit{informed} about widespread misconceptions—for example, a false belief held by half the population—but not reproduce or assert this misconception in half of its responses. Thus, discriminating between the model’s broad repertoire of \emph{knowledge and skills} and specifying its \emph{targeted responses and behaviors} is essential for ensuring the development of AI systems that are safe, effective, and reliably controllable \citep{ouyang2022training}. Most current strategies guide LMs towards human-valued outputs through reinforcement learning (RL), yet we will demonstrate that the RL objective frequently adopted can, in fact, be precisely solved via a straightforward binary cross-entropy loss, offering a more direct and streamlined approach to preference learning.

In general, current approaches adapt a language model to exhibit preferred behaviors using carefully constructed datasets of human preferences that exemplify responses deemed safe and useful by humans. This preference learning step is applied subsequent to a large-scale, unsupervised pre-training phase on textual data. Although the most direct route to preference alignment is to employ supervised fine-tuning based on exemplary human demonstrations, the reinforcement learning from human (or AI) feedback paradigm—RLHF or RLAIF~\citep{christiano2017deep,bai2022constitutional}—has become the leading methodology. RLHF first fits a reward model to human-preference-annotated data, and then applies reinforcement learning to adjust the LM policy, promoting generation of responses with higher reward values, while attempting to avoid significant deviation from the pre-trained model. Although RLHF methods achieve impressive results in tasks such as dialogue and code generation, they impose considerable complexity compared to simple supervised approaches, typically requiring the training of multiple language models and integrating in-the-loop sampling from the model policy during training, all of which substantially increase computational requirements.

This work demonstrates a method for optimizing language models to align with human preferences directly, bypassing the need for explicit reward modeling or reinforcement learning frameworks. We introduce \textit{{Direct Preference Optimization} (DPO)}, a technique that achieves the same objective as traditional RLHF approaches—reward maximization subject to a KL-divergence regularization—while offering a more straightforward implementation and training procedure. Fundamentally, DPO works by adjusting the model to increase the log-likelihood ratio between preferred and non-preferred outputs, while employing adaptive, instance-level importance weights to avoid model degradation observed with simplistic probability ratio objectives. In keeping with prior art, DPO leverages a theoretical model of preference, such as the Bradley-Terry model (\cite{bradley1952rankanalysis}), which provides a measure of the congruence between a reward function and actual observed preferences. However, instead of using the preference model to construct a loss for reward model training—followed by separate policy optimization on this learned reward—DPO applies a change-of-variables approach to define the preference loss directly with respect to the policy. Utilizing a dataset composed of human judgments over responses, DPO is able to update a policy via a basic binary cross entropy loss, leading to the optimal policy for an implicitly defined reward function fitted to the collected preference data.

The central contribution of our work is the Direct Preference Optimization (DPO) algorithm, a reinforcement learning-free approach for preference-driven language model training. Experimental results indicate that DPO achieves competitive, if not superior, performance compared to current alternatives, such as PPO-based RLHF, in domains including sentiment modification, text summarization, and conversation, across models scaling up to 6B parameters.

Large-scale self-supervised language models demonstrate the ability to perform certain tasks in a zero-shot fashion \citep{radford2019language} or with a minimal number of prompt examples \citep{gpt3,megatron,chowdhery2022palm}. Nonetheless, substantial enhancements in downstream task performance and alignment with user expectations are achieved through fine-tuning on curated datasets comprising instructions and human-generated outputs \citep{mishra-etal-2022-cross,sanh2022multitask,chung2022scaling,thoppilan2022lamda}. This approach, termed `instruction-tuning,' improves the capacity of LLMs to handle instructions beyond the original fine-tuning set, thereby generally enhancing model utility \citep{chung2022scaling}. While instruction tuning has produced notable advances, collecting \textit{relative} human preference judgments regarding output quality tends to be more scalable than gathering expert-level demonstrations. Consequently, more recent studies have focused on refining LLMs using datasets reflecting human preferences, which has led to advancements in areas like translation \citep{kreutzer-etal-2018-reliability}, summarization \citep{stiennon2022learning,ziegler2020finetuning}, narrative generation \citep{ziegler2020finetuning}, and following instructions \citep{ouyang2022training,ramamurthy2023is}. The common procedure involves first training a neural network-based reward model to align with preference data—frequently employing frameworks such as the Bradley-Terry model \citep{bradley1952rankanalysis}—and subsequently adapting a language model to optimize this reward through reinforcement learning techniques, such as REINFORCE \citep{williams1992reinforce}, proximal policy optimization (PPO; \cite{schulman2017proximal}), or their variants \citep{ramamurthy2023is}. A related research direction utilizes LLMs already tuned to follow instructions and fine-tuned with human feedback to create synthetic preference datasets focused on specific properties like safety or harmlessness \citep{bai2022constitutional}, relying only on minimal human input such as text-based rubrics to guide LLM annotations. These strategies represent a synthesis of two research traditions: one on reinforcement learning-based language model training across various tasks~\citep{Ranzato2015SequenceLT,paulus2018a,wu2018learning}, and the other on general techniques for preference-based learning from human feedback \citep{christiano2017deep,kupcsik2018learning}. Although leveraging human preferences via relative judgments is an attractive approach, applying reinforcement learning to fine-tune large language models remains highly complex in practice; the current work introduces a methodologically sound alternative for optimizing relative preferences without the need for RL.

Beyond the scope of language tasks, the study of learning policies based on preferences has garnered attention in both bandit and reinforcement learning paradigms, with a variety of methods emerging in the literature. When preferences or rankings over actions replace explicit reward signals, this setting is known as contextual dueling bandits (CDB; \cite{yue2012karmed,dudik2015contextual}). In these cases, where absolute reward values are not present, CDB theory introduces the concept of a \textit{von Neumann winner}, defined as a policy that achieves an expected win rate of no less than 50\% when pitted against any other policy \citep{dudik2015contextual}. Distinct from the RLHF scenario, CDB frameworks rely on receiving preference information in an online fashion, whereas learning from human preferences typically involves leveraging a pre-collected, finite dataset of preference-labeled action pairs \citep{yan2022human}. A related line of work, known as \textit{preference-based RL} (PbRL), employs binary preferences derived from an unobserved `scoring' function in lieu of explicit reward feedback \citep{BusaFekete2014,ruiz2023dueling}. Numerous strategies for PbRL have been proposed, often allowing for the utilization of off-policy preference data; these methods generally require first inferring the hidden scoring or reward function, and then optimizing with respect to this estimate \citep{jain2013learning,BusaFekete2014,christiano2017deep,sadigh2017active,kupcsik2018learning}. In contrast, our approach bypasses explicit reward model estimation, introducing a one-step procedure that directly trains a policy to accommodate given preferences.

\section{Preliminaries}\label{section:prelims}

We summarize the RLHF pipeline described by \citeauthor{ziegler2020finetuning} and subsequently in \citep{stiennon2022learning, bai2022training, ouyang2022training}. This framework consists of three principal stages: 1) supervised fine-tuning (SFT); 2) the preference elicitation and reward modeling stage; and 3) reinforcement learning optimization.

\textbf{SFT}: The RLHF process begins with supervised fine-tuning, in which a pre-trained language model is further trained on curated, high-quality datasets tailored to target applications (such as dialogue or summarization). This step yields the model $\pi^\text{SFT}$.

\textbf{Reward Modelling Phase}: The next stage involves using the SFT model to respond to prompts $x$, producing answer pairs $(y_1, y_2)\sim \pi^\text{SFT}(y \mid x)$. Human annotators are then tasked with indicating which of the two answers they prefer, written as $y_w\succ y_l \mid x$, where $y_w$ and $y_l$ represent the chosen and non-chosen completions from the pair. The observed preferences are assumed to arise from some unobserved reward function $r^*(y, x)$, which is not directly available. Several preference modeling techniques are applicable at this stage, with the Bradley-Terry (BT) \cite{bradley1952rankanalysis} model being especially prevalent (more complex Plackett-Luce models \citep{plackett1975analysis, luce2012individual} can also be used, especially with multi-way rankings). According to the BT formulation, the distribution $p^*$ over human preferences can be represented as follows:

Given a fixed comparison dataset $\mathcal{D} = \bigl\{x^{(i)}, y_w^{(i)}, y_l^{(i)}\bigr\}_{i=1}^N$ drawn according to $p^*$, we introduce a parameterized reward function $r_{\phi}(x, y)$, whose parameters are learned through maximum likelihood estimation. Treating this as a binary classification task, the associated negative log-likelihood loss function is given by:

\begin{equation}\label{eq:reward_model}
    \mathcal{L}_R(r_{\phi}, \mathcal{D}) = -\mathbb{E}_{(x, y_w, y_l)\sim \mathcal{D}}\left[\log \sigma\left(r_{\phi}(x, y_w)- r_{\phi}(x, y_l)\right)\right],
\end{equation}

where $\sigma$ is the standard logistic (sigmoid) function. For language models, $r_{\phi}(x, y)$ commonly reuses the architecture of the SFT model $\pi^\text{SFT}(y \mid x)$, extending it with an additional linear projection applied to the last transformer layer, which outputs a scalar score corresponding to the reward \cite{ziegler2020finetuning}. In order to limit the variance of the learned reward, previous studies standardize the outputs so that $\mathbb{E}_{x,y\sim \mathcal{D}}[r_\phi(x, y)] = 0$ across the dataset for any $x$.

\textbf{RL Fine-Tuning Phase}: Once the reward model has been trained, it is employed during reinforcement learning to guide the language model's updates. Following established work~\citep{jaques2017sequence, jaques2020human}, this process is formalized by the following objective:

\begin{equation}\label{eq:RL}
\max_{\pi_{\theta}}  \mathbb{E}_{x\sim \mathcal{D}, y\sim \pi_{\theta}(y \mid x)}\left[r_{\phi}(x, y)\right] - \beta\mathbb{D}_{\textrm{KL}}\left[\pi_{\theta}(y\mid x)\, \Vert\, \pi_\text{ref}(y\mid x)\right],
\end{equation}

where $\beta$ modulates the distance from the initial policy $\pi_\text{ref}$ (specifically, $\pi^\text{SFT}$). The language model policy $\pi_{\theta}$ is typically initialized from $\pi^\text{SFT}$. The inclusion of the KL-divergence constraint serves to maintain alignment with the region of policy space where the reward model is reliable, support output diversity, and mitigate the risk of collapse into repetitive, high-reward outputs. Because text generation operates over discrete tokens, this objective is not differentiable and thus necessitates reinforcement learning approaches. The typical strategy \citep{ziegler2020finetuning, stiennon2022learning, bai2022training, ouyang2022training} modifies the reward to the form ${r(x, y) = r_{\phi}(x, y) -\beta (\log \pi_{\theta}(y\mid x) - \log \pi_\text{ref}(y\mid x))}$, optimizing it with PPO \cite{schulman2017proximal}.

\section{Direct Preference Optimization}\label{sec:DPO}

Recognizing the practical difficulties of deploying reinforcement learning algorithms for large-scale models, such as those used in language model fine-tuning, our aim is to construct a streamlined method for policy optimization based on direct preference data. Rather than separately training a reward model and then employing reinforcement learning as in conventional RLHF approaches, our method leverages a specific reward model parameterization that admits a closed-form expression for the corresponding optimal policy—eliminating the need for a traditional RL loop. Our central insight lies in establishing an analytical correspondence between reward models and their optimal policies, which allows us to transform loss objectives from the reward parameter space into the policy space. By applying this change-of-variables, we can bypass explicit reward model training while still optimizing with respect to established models of human preference, such as the Bradley-Terry framework. Effectively, our approach unifies the policy network and the (implicit) reward, enabling the language model itself to encode both functionalities.

\textbf{Deriving the DPO objective.} We begin from the general reinforcement learning objective given in Eq.~\ref{eq:RL}, which employs a reward function $r$, following conventions established in prior research~\citep{peters2007reinforcement, peng2019advantage, korbak2022reinforcement, go2023aligning}. It can be readily demonstrated that the solution which optimally maximizes reward while constrained by a KL divergence, as in Eq.~\ref{eq:RL}, adopts the following form:

\begin{equation}\label{eq:op_policy}
    \pi_r(y\mid x) = \frac{1}{Z(x)}\pi_\text{ref}(y\mid x)\exp\left(\frac{1}{\beta}r(x, y)\right),
\end{equation}

where the partition function $Z(x) =\sum_{y}\pi_\text{ref}(y\mid x)\exp\left(\frac{1}{\beta}r(x, y)\right)$. A full derivation is included in Appendix \ref{app:derivation1}. Despite estimating the reward function using an MLE approximation $r_\phi$ for the true reward $r^*$, computing $Z(x)$ remains computationally costly~\citep{korbak2022reinforcement, go2023aligning}, which reduces the tractability of this direct approach. However, it is possible to algebraically transform Eq.~\ref{eq:op_policy} to rewrite the reward in terms of the optimal policy $\pi_r$, the reference policy $\pi_\text{ref}$, and $Z(\cdot)$. Specifically, by taking logarithms of both sides and rearranging, we find:

\begin{equation}\label{eq:main_eq}
    r(x,y) =\beta \log \frac{\pi_r(y\mid x)}{\pi_\text{ref}(y\mid x)} + \beta \log Z(x).
\end{equation}

This transformation also applies to the ground-truth reward $r^*$ and the optimal policy $\pi^*$. Importantly, in the Bradley-Terry model, only the difference between the rewards for two different completions is relevant; that is, ${p^*(y_1 \succ y_2 \mid x) = \sigma(r^*(x, y_1) - r^*(x, y_2))}$. Substituting the expression for $r^*(x,y)$ from Eq.~\ref{eq:main_eq} into the preference model of Eq.~\ref{eq:bradley-terry}, the partition function cancels out, leaving the preference probability dependent only on the optimal policy $\pi^*$ and the reference $\pi_\text{ref}$. Thus, under the Bradley-Terry preference model, the optimal RLHF policy $\pi^*$ fulfills:

\begin{equation}\label{eq:objective}
    p^*(y_1\succ y_2 \mid x)=\frac{1}{1 + \exp\left(\beta \log \frac{\pi^*(y_2\mid x)}{\pi_\text{ref}(y_2\mid x)} - \beta \log \frac{\pi^*(y_1\mid x)}{\pi_\text{ref}(y_1\mid x)}\right)}
\end{equation}

A detailed derivation is provided in Appendix~\ref{app:derivation2}. Although Eq.~\ref{eq:objective} is based on the Bradley-Terry model, parallel derivations for the broader Plackett-Luce models~\citep{plackett1975analysis, luce2012individual} can also be obtained, as shown in Appendix~\ref{app:plackett_luce_models}.

With the probability of observed human preferences now formulated as a function of the optimal policy rather than the reward model, it is possible to develop a maximum likelihood estimation objective for any parameterized policy $\pi_\theta$. Mirroring the structure of the reward-based modeling objective (see Eq.~\ref{eq:reward_model}), our objective for training the policy takes the form:

\begin{equation}\label{eq:optimum_model}
    \mathcal{L}_\text{DPO}(\pi_{\theta}; \pi_\text{ref}) = -\mathbb{E}_{(x, y_w, y_l)\sim \mathcal{D}}\left[\log \sigma \left(\beta \log \frac{\pi_{\theta}(y_w\mid x)}{\pi_\text{ref}(y_w\mid x)} - \beta \log \frac{\pi_{\theta}(

In this formulation, we learn an implicit reward through an alternative parametrization in which the optimal policy corresponds directly to $\pi_\theta$. Additionally, since our approach amounts to training a reparametrized Bradley-Terry model, it inherits various desirable theoretical guarantees, such as consistency provided the preference data distribution satisfies certain conditions \cite{bong2022generalized}. We further elaborate on the theoretical aspects of DPO and its comparison to related methods in Section~\ref{sec:theory}.

\textbf{Interpretation of the DPO update.} To develop a more mechanistic understanding of DPO, it is instructive to examine the gradient of its objective, $\mathcal{L}_\text{DPO}$. The gradient with respect to the model parameters $\theta$ is given by:

\begin{multline*}\label{eq:gradient}
    \nabla_\theta \mathcal{L}_\text{DPO}(\pi_\theta;\pi_\text{ref}) = \\ -\beta\mathbb{E}_{(x, y_w, y_l) \sim \mathcal{D}} \bigg[\underbrace{\sigma(\hat{r}_\theta(x, y_l) - \hat{r}_\theta (x, y_w))}_\text{larger magnitude if reward ordering is incorrect}\bigg[\underbrace{\nabla_\theta\log \pi(y_w \mid x)}_\text{promote $y_w$} - \underbrace{\nabla_\theta\log\pi(y_l \mid x)}_\text{demote $y_l$}\bigg]\bigg],
\end{multline*}

where the implicit reward, $\hat{r}_\theta(x, y) = \beta \log \frac{\pi_\theta(y \mid x)}{\pi_\text{ref}(y \mid x)}$, is determined by the ratio of the model $\pi_\theta$ to the reference $\pi_\text{ref}$ (see Section~\ref{sec:theory} for further details). Conceptually, this gradient update increases the probability assigned to preferred completions $y_w$ while decreasing it for less preferred completions $y_l$. The weighting for each example is determined by the extent to which the implicit reward model $\hat{r}_\theta$ erroneously ranks $y_l$ above $y_w$, modulated by $\beta$, thereby reflecting the model's current misclassification and the influence of the KL regularization. Empirically, this weighting proves essential; omitting it (i.e., using a variant of the method without this coefficient) can lead to undesirable behaviors such as model degeneration, as evidenced in Appendix Table~\ref{tab:unlikelihood_generations}.

\textbf{Overview of DPO.} The standard DPO workflow typically involves: 1) For each prompt $x$, generating two responses $y_1, y_2 \sim \pi_\text{ref}(\cdot \mid x)$, and annotating them with human preferences to produce the preference dataset $\mathcal{D} = \{x^{(i)}, y_w^{(i)}, y_l^{(i)}\}_{i=1}^N$; and 2) Training the language model $\pi_\theta$ to minimize the DPO loss $\mathcal{L}_\text{DPO}$, conditioned on the specified $\pi_\text{ref}$, the constructed dataset $\mathcal{D}$, and the chosen hyperparameter $\beta$.  
In real-world scenarios, practitioners prefer to utilize already existing public preference datasets, rather than creating new samples and collecting fresh human annotations. Since these datasets are generated under $\pi^\text{SFT}$, we set $\pi_\text{ref} = \pi^\text{SFT}$ where possible. If $\pi^\text{SFT}$ is absent, we instead estimate $\pi_\text{ref}$ by maximizing the likelihood over the preferred completions ${(x, y_w)}$: that is, ${\pi_\text{ref} = \argmax_{\pi}\mathbb{E}_{x, y_w \sim \mathcal{D}}\left[\log \pi(y_w \mid x)\right]}$. This initialization aims to reduce mismatch between the inaccessible true reference distribution and the $\pi_\text{ref}$ used during DPO training. Additional implementation and hyperparameter information is available in Appendix~\ref{app:implementation}.

\section{Theoretical Analysis of DPO}
This section seeks to deepen understanding of the DPO method, establish its theoretical foundation, and connect its benefits to limitations observed in actor-critic methods utilized for RLHF, such as PPO~\cite{schulman2017proximal}.

\label{sec:theory}

\subsection{Your Language Model Is Secretly a Reward Model}
DPO effectively merges reward inference and policy optimization by leveraging a single maximum likelihood framework, foregoing the need for separate explicit reward modeling and RL-based policy updates. The objective in Eq. \ref{eq:main_eq} can be interpreted as corresponding to a Bradley-Terry model with the reward function parameterized as $r^*(x, y) = \beta \log\frac{\pi^*_\theta(y \mid x)}{\pi_\text{ref}(y \mid x)}$. We fit our parametric model $\pi_\theta$ directly, paralleling the optimization performed on the reward model in Eq. \ref{eq:reward_model} after an appropriate change of variables. In what follows, we develop the theoretical justification for this parameterization, demonstrate that it imposes no inherent limitations on the representational power of the reward model class, and show that this approach achieves the optimal policy. To do so, we start by introducing an equivalence relation for reward functions.

\begin{definition}
Two reward functions $r(x, y)$ and $r'(x, y)$ are said to be equivalent if and only if ${r(x, y)-r'(x, y) = f(x)}$ for some function $f$.
\end{definition}

This relationship clearly satisfies the conditions for an equivalence relation, thereby dividing the set of all reward functions into equivalence classes. The following lemmas capture important properties of these classes:

\begin{lemma}\label{lemma:same_prefrence}
Within the Plackett-Luce framework—including the Bradley-Terry model—any two reward functions belonging to the same equivalence class generate identical preference distributions.
\end{lemma}

\begin{lemma}\label{lemma:same_policy}
For the constrained RL formulation, all reward functions within a given equivalence class result in the same optimal policy.
\end{lemma}

The arguments supporting these results are direct and are provided in Appendix~\ref{app:lemma1}. The first lemma highlights a well-known identifiability problem inherent to the Plackett-Luce model family \cite{plackett1975analysis}. This ambiguity necessitates further constraints to enable identifiability and to obtain meaningful MLE solutions, as noted in Eq.~\ref{eq:reward_model} and discussed in \cite{bong2022generalized}. The second lemma ensures that any member of a reward equivalence class yields the same optimal policy, thus justifying our focus on recovering an arbitrary representative from this class as our ultimate goal. The theorem below, proved in Appendix~\ref{app:thm1}, formalizes the representation of these classes:

\begin{theorem}\label{thm:main}
Provided that mild conditions are satisfied, every reward equivalence class consistent with the Plackett-Luce (or specifically, the Bradley-Terry) model can be expressed through the reparameterization ${r(x, y) = \beta \log \frac{\pi(y\mid x)}{\pi_\text{ref}(y\mid x)}}$, where $\pi(y\mid x)$ denotes a model and $\pi_\text{ref}(y\mid x)$ is a specified reference model.
\end{theorem}

\begin{sproof}
Let $r(x, y)$ denote any given reward function that induces the optimal policy $\pi_r(y \mid x)$ as characterized by Eq.~\ref{eq:op_policy}. We aim to demonstrate that any reward function within the equivalence class of $r$ admits the stated reparameterized form. To this end, we define the following projection:
\begin{equation}
    f(r; \pi_\text{ref}, \beta)(x, y) = r(x, y) - \beta\log\sum_{y}\pi_\text{ref}(y\mid x)\exp\left(\frac{1}{\beta}r(x, y)\right)
\end{equation}
This operator $f$ rescales the reward function by subtracting the log of its partition function (which depends only on $x$), yielding a new reward function that remains in the same equivalence class as $r(x, y)$. Substituting the expression from Eq.~\ref{eq:main_eq} (which universally holds for any reward), it follows that $f(r; \pi_\text{ref}, \beta)(x, y) = \beta \log \frac{\pi_r(y\mid x)}{\pi_\text{ref}(y\mid x)}$. Therefore, $f$ produces a reward in the equivalence class of $r$ with the prescribed structure, ensuring no loss of generality when adopting this reparameterization for our reward model.
\end{sproof}

An alternative interpretation of Theorem~\ref{thm:main} is that it identifies the specific reward function chosen by the DPO reparameterization within each equivalence class—namely, the one fulfilling:

\begin{equation}\label{eq:lag_p}
     \sum_{y}\underbrace{\pi_\text{ref}(y\mid x)\exp\left(\frac{1}{\beta}r(x, y)\right)}_{=\pi(y\mid x)\text{, using Thm.~\ref{thm:main} reparam.}} = 1,
\end{equation}

This condition ensures that $\pi(y\mid x)$ forms a legitimate probability distribution, where all probabilities are non-negative and their total sums to 1. If we reconsider Eq.~\ref{eq:lag_p} from the perspective of Eq.~\ref{eq:op_policy}, we observe that this equation represents the partition function associated with the optimal policy determined by $r(x, y)$. The essential insight behind the DPO method is its ability to enforce particular constraints on the inherently under-determined Plackett-Luce preference models (and, by extension, the Bradley-Terry family), thereby retaining the set of achievable reward models. Crucially, this makes the optimal policy characterized in Eq.~\ref{eq:op_policy} explicitly computable in closed form for any prompt $x$.

\subsection{Instability of Actor-Critic Algorithms} Our framework further allows us to analyze sources of instability found in typical actor-critic approaches in RLHF, such as PPO. Concentrating on the RL fine-tuning component as detailed in Section \ref{section:prelims}, and situating our analysis within the control as inference paradigm \cite{levine2018reinforcement} for the constrained RL formulation in \ref{eq:RL}, we parameterize the model as $\pi_{\theta}(y\mid x)$. The objective then becomes minimizing $\mathbb{D}_{\text{KL}}[\pi_{\theta}(y|x) \mid \mid \pi^*(y\mid x)]$, where the reference $\pi^*$ denotes the optimal policy derived in Eq. \ref{eq:optimum_model} via reward function $r_{\phi}(y, x)$. After algebraic manipulation, the optimization target is given by:

\begin{equation}\label{eq:AC}
    \max_{\pi_{\theta}}\mathbb{E}_{\pi_{\theta}(y\mid x)}\bigg[\underbrace{r_{\phi}(x, y) -\beta\log\sum_{y}\pi_\text{ref}(y\mid x)\exp\left(\frac{1}{\beta}r_{\phi}(x, y)\right)}_{f(r_{\phi}, \pi_\text{ref}, \beta)} - \underbrace{\beta\log\frac{\pi_{\theta}(y\mid x)}{\pi_\text{ref}(y\mid x)}}_{\text{KL}}\bigg]
\end{equation}

This objective has also been the focus of optimization in previous studies 
\citep{ziegler2020finetuning, stiennon2022learning, bai2022training, ouyang2022training}, where the DPO-equivalent reward corresponding to the reward family $r_{\phi}$ is utilized. Within this context, the normalization component of $f(r_{\phi}, \pi_\text{ref}, \beta)$ can be viewed as the soft value function associated with the reference policy $\pi_\text{ref}$. Although this normalization does not alter the optimum of the objective, omitting it can lead to high-variance policy gradients, which in turn may destabilize the learning process. A possible remedy involves approximating the normalization with a trainable value function, but this introduces its own optimization challenges. Alternatively, earlier works have addressed normalization by employing a human completion baseline, effectively using a single-sample Monte-Carlo estimate to approximate the normalization term. By contrast, the DPO reparameterization directly provides a reward formulation that eliminates the need for such baselines.

\section{Experiments}
Here, we conduct empirical analyses to assess DPO's capacity to train policies purely from preference feedback. We begin by examining a controlled text-generation scenario to investigate how effectively DPO balances reward maximization and KL-divergence regularization relative to standard preference-based methods like PPO. Subsequently, we extend our evaluation to more complex RLHF tasks and larger model scales, considering challenges such as summarization and dialogue generation. Our findings indicate that, even with minimal hyperparameter adjustment, DPO generally matches or surpasses established baselines, including RLHF with PPO and the approach of selecting the highest-reward trajectory among $N$ samples. We outline our experimental methodology below, with further details provided in Appendix~\ref{app:exp_details}.

\textbf{Tasks.} Our study examines three distinct open-ended text generation tasks. Across all tasks, algorithms derive a policy using a preference dataset $\mathcal{D}=\bigl\{x^{(i)}, y_w^{(i)}, y_l^{(i)}\bigr\}_{i=1}^N$. In the context of \textbf{controlled sentiment generation}, $x$ represents the start of a movie review sourced from the IMDb dataset \cite{maas-EtAl:2011:ACL-HLT2011}, and the goal of the policy is to produce $y$ exhibiting positive sentiment. To allow for a systematic assessment, we synthesize preference pairs by generating continuations and leveraging a pre-trained sentiment classifier to select pairs such that $p(\text{positive}\mid x,y_w)>p(\text{positive}\mid x,y_l)$. For SFT, GPT-2-large is fine-tuned to convergence on the training portion of the IMDB review data (for additional methodology, see App~\ref{app:sentiment_details}). Within the \textbf{summarization} task, $x$ corresponds to a Reddit forum post, and the policy’s objective is to output a summary $y$ that captures the primary points. Building on previous research, we utilize the Reddit TL;DR summarization dataset \citep{volske-etal-2017-tl} along with preference data collected by \citeauthor{stiennon2022learning}. SFT is conducted on a model further trained using human-generated forum summaries,\footnote{\url{https://huggingface.co/CarperAI/openai_summarize_tldr_sft}} and RLHF is implemented using the TRLX framework \citep{leandro_von_werra_2023_7790115}. The human preference dataset originates from \citeauthor{stiennon2022learning}, who selected samples generated by a comparable SFT model. In the \textbf{single-turn dialogue} task, $x$ is a prompt from a user, which might range from inquiries on astrophysics to requests for relationship guidance. Here, the policy is tasked with composing an informative and engaging answer $y$; we employ the Anthropic Helpful and Harmless dialogue dataset \citep{bai2022training}, encompassing 170,000 human–assistant dialogues. Each interaction concludes with two responses created by a large, unspecified language model, one of which is annotated as preferred by a human rater. Unlike prior tasks, a pre-trained SFT model is not available in this context, so we fine-tune a publicly available language model exclusively on the preferred responses to establish the SFT baseline.

\textbf{Evaluation.} Our experimental evaluation utilizes two complementary strategies. For a controlled investigation into how well each algorithm addresses the constrained reward maximization problem, we employ the controlled sentiment generation scenario, where we assess each method based on its performance frontier—measured by achieved reward versus KL-divergence relative to the reference policy. The availability of a ground-truth reward function (provided by a sentiment classifier) enables precise computation of this trade-off. Conversely, when the ground-truth reward function is inaccessible, as is typically the case in real-world situations, we turn to empirical win rates against a baseline policy, using GPT-4 as an automated stand-in for human judgments regarding summary quality and helpfulness in both the summarization and single-turn dialogue domains. Here, baseline policies consist of either the reference summaries included in the test data for summarization or the preferred response chosen within the dialogue test set. While prior work indicates large language models can surpass established metrics in evaluation fidelity \citep{Chen2023ExploringTU}, we supplement our protocol with a human annotation study to further validate the use of GPT-4 as an evaluator, as discussed in Sec.~\ref{sec:human-judgments}. Our findings show that human agreement with GPT-4’s assessments is high, often matching or even exceeding agreement between human annotators.

\textbf{Methods.} Beyond DPO, we examine a variety of established training approaches for aligning language models with human preferences. Our simplest baselines utilize zero-shot prompting with \textbf{GPT-J} \citep{gpt-j} for summarization and 2-shot prompting with \textbf{Pythia-2.8B} \citep{biderman2023pythia} for dialogue. We also consider the standard \textbf{SFT} model, alongside \textbf{Preferred-FT}—which is trained by supervised fine-tuning on selected completions $y_w$ obtained either from the SFT model (for sentiment control and summarization tasks) or from a general-purpose language model (for dialogue). An additional supervised-style baseline is the \textbf{Unlikelihood} approach~\citep{welleck2019neural}, which optimizes the policy by encouraging the model to increase the likelihood of $y_w$ while decreasing that of $y_l$, with the negative component weighted by a tunable parameter $\alpha\in[0,1]$. Reinforcement learning methods are also explored, namely \textbf{PPO} \citep{schulman2017proximal}—using a reward function inferred from preference data—and \textbf{PPO-GT}, an oracle model which benefits from direct access to the ground-truth reward during controlled sentiment experiments. In these sentiment experiments, we implement two versions of PPO-GT: an out-of-the-box solution \cite{leandro_von_werra_2023_7790115} and a modified variant that incorporates reward normalization and optimized hyperparameters to enhance outcomes; these modifications are likewise applied to standard PPO with learned rewards. Lastly, we introduce the \textbf{Best of $N$} baseline, which entails generating $N$ candidate outputs via the SFT model (or Preferred-FT in dialogue settings) and selecting the response with the highest score according to a reward model trained on preference data. Despite its strong empirical performance, this approach is computationally burdensome, as it requires generating $N$ samples for every input during evaluation.

\subsection{How well can DPO optimize the RLHF objective?}

The standard RLHF algorithms employ a reward maximization objective constrained by the KL divergence, balancing the need to obtain high rewards while ensuring that the learned policy remains close to a reference policy. Thus, in evaluating different algorithms, both the magnitude of reward obtained and the extent of KL divergence from the reference must be considered; higher rewards accompanied by a significant increase in KL may not represent a favorable outcome. Figure~\ref{fig:frontier-tldr-main} illustrates the relationship between reward and KL divergence for several methods within the sentiment analysis context. For each approach, we conduct multiple independent training sessions, each utilizing a unique value for the policy conservativeness parameter (target KL $\in\{3,6,9,12\}$ for PPO, $\beta \in \{0.05,0.1,1,5\}$, $\alpha\in\{0.05,0.1,0.5,1\}$ for unlikelihood, as well as random seeds for preferred-FT), amounting to a total of 22 different experimental runs. At intervals of 100 training steps, policies are assessed using a designated test set of prompts, and we calculate both the mean reward based on the actual reward function and the average KL across sequences\footnote{Here, KL refers to the total of per-timestep KL-divergence values.} with respect to the reference policy, specifically $\text{KL}\left(\pi\mid \mid \pi_\text{ref}\right)$. The findings demonstrate that DPO yields a notably superior frontier, reaching greater rewards without incurring excessive KL divergence. This outcome is striking for several reasons: while DPO and PPO target the same optimization problem, DPO demonstrates substantially improved efficiency, with its reward/KL performance strictly surpassing that of PPO. Furthermore, DPO achieves an improved reward-KL tradeoff even in comparison to PPO-GT, where PPO is supplied with exact reward signals.

\subsection{Can DPO scale to real preference datasets?}
\label{sec:dpo-real-datasets}
We proceed to assess the effectiveness of DPO fine-tuning on real-world tasks, specifically in the domains of summarization and single-turn dialogue. In the context of summarization, traditional automatic metrics like ROUGE often display a weak correlation with human preferences~\citep{stiennon2022learning}. Earlier research has demonstrated that training language models with PPO guided by human feedback can yield more compelling summaries. Our experimental setup involves generating outputs from various methods on the test split of the TL;DR summarization dataset and then calculating the mean win rate compared to the reference summaries provided in the test set. For each method, completions are drawn across a temperature range of 0.0 to 1.0, with the resulting win rates visualized in Figure~\ref{fig:frontier-tldr-main} (right). All methods, including DPO, PPO, and Preferred-FT, use the same GPT-J SFT base model for fine-tuning\footnote{\url{https://huggingface.co/CarperAI/openai_summarize_tldr_sft}}. Notably, DPO achieves a win rate near 61\% at a temperature setting of 0.0, surpassing PPO, which attains around 57\% under its optimal temperature condition. Furthermore, DPO outperforms the best of $N$ baseline by achieving a higher peak win rate. It is important to mention that no substantial tuning of DPO's $\beta$ hyperparameter was conducted, indicating these outcomes may not fully reflect DPO's capabilities. Additionally, DPO demonstrates superior robustness to temperature fluctuations compared to PPO, whose performance can diminish to the level of the baseline GPT-J at higher temperatures. The Preferred-FT approach, meanwhile, does not yield marked improvements over the SFT baseline. Finally, in direct human assessments detailed in Section~\ref{sec:human-judgments}, DPO samples generated at a temperature of 0.25 were favored 58\% of the time compared to PPO samples produced at a temperature of 0.

For single-turn dialogue evaluation, we assess various methods using the subset of the Anthropic HH dataset \citep{bai2022training} test split containing a single round of human-assistant exchange. To determine win rates, GPT-4 evaluations use the test set's preferred completions as the benchmark across all methods. Since a standard SFT model is not established for this setting, we begin with a pre-trained Pythia-2.8B model, employ Preferred-FT to create a reference model trained on the selected completions so that its outputs remain within the model's distribution, and subsequently apply DPO for further training. Our comparisons include the strongest result from 128 Preferred-FT completions (noting that the Best of $N$ method yields diminishing returns past 128 samples, as discussed in Appendix Figure~\ref{fig:best-of-n}) and a 2-shot prompt variant of the original Pythia-2.8B model. Results indicate that, for optimal temperature settings per method, DPO consistently matches or outperforms alternatives. Additionally, we tested an RLHF model utilizing PPO, trained on the Anthropic HH dataset\footnote{\url{https://huggingface.co/reciprocate/ppo_hh_pythia-6B}} and released by a reputable source\footnote{\url{https://github.com/CarperAI/trlx/tree/main/examples/hh}}. However, even after experimenting with various prompts and sampling temperatures, we observed no improvement over the unmodified Pythia-2.8B model. Given our findings from TL;DR and the alignment of reward optimization objectives between both approaches, we treat Best of 128 as an approximate indicator for PPO-level performance. In summary, DPO emerges as the only computationally efficient technique that surpasses the test set preferred completions from the Anthropic HH dataset, and it achieves performance on par with, or superior to, the computationally intensive Best of 128 approach. Furthermore, as illustrated in Figure~\ref{fig:dialogue-main}, DPO rapidly attains peak performance.

\subsection{Generalization to a new input distribution}

To assess how well PPO and DPO cope with distributional changes, we test the summarization policies trained on Reddit TL;DR data against news articles from the CNN/DailyMail dataset's test split \citep{nallapati-etal-2016-abstractive}. For this evaluation, we retain the best sampling temperatures identified on the TL;DR data (0 and 0.25). Table~\ref{tab:ood} summarizes our findings. To measure performance, we determine the win rate of GPT-4 when choosing between model-generated summaries and ground-truth references, applying the same GPT-4 (C) evaluation prompt used in the Reddit experiments, with "forum post" replaced by "news article". Across this new domain, DPO maintains a considerable advantage over PPO. These results suggest that DPO-trained policies can generalize at least as effectively as those from PPO, despite the absence of additional unlabeled Reddit TL;DR data during DPO training.

\subsection{Validating GPT-4 judgments with human judgments}
\label{sec:human-judgments}

We carry out a human evaluation study to determine how closely GPT-4’s assessments align with human preferences, focusing on the TL;DR summarization task and two different GPT-4 prompts. The \textbf{GPT-4 (S)} (simple) prompt asks which summary more effectively conveys the key information from the source. The \textbf{GPT-4 (C)} (concise) prompt inquires about both informativeness and conciseness, the latter included because we observe that with the simple prompt, GPT-4 prefers summaries that are longer and more repetitive than human raters do. The full prompt text is included in Appendix~\ref{app:prompts}. Three sets of model comparisons are performed—featuring the highest performing model (DPO, temp. 0.25), the lowest (PPO, temp. 1.0), and an intermediate system (SFT, temp. 0.25)—against greedily-sampled PPO, which performed best at its optimal temperature. This approach captures a range of output qualities. The results indicate that, under both prompts, GPT-4 agrees with human judgments at rates comparable to inter-human agreement, implying that GPT-4 can function as a suitable stand-in for human raters (due to a small pool of raters, only the DPO and PPO-1 cases receive multiple independent human judgments). Overall, the \textbf{GPT-4 (C)} prompt aligns most closely with human preferences, leading us to adopt it as the principal evaluation measure in Section~\ref{sec:dpo-real-datasets}. Appendix~\ref{app:human-study} provides further methodological details, including information about the evaluation interface and human participants.

\section{Discussion}
Preference-based learning presents a robust and scalable approach for developing competent and value-aligned language models. In this work, we introduced DPO, a straightforward methodology for utilizing preferences in language model training that avoids the complexities of reinforcement learning. Instead of reframing preference learning to fit into conventional RL frameworks for compatibility with standard algorithms, DPO establishes a correspondence between language model policies and reward mechanisms, thereby allowing direct optimization towards human preferences through a simple cross-entropy objective, eliminating the need for reinforcement learning and preserving general applicability. DPO demonstrates performance on par with, or superior to, established RLHF techniques—such as those employing PPO—while requiring minimal hyperparameter adjustment. As a result, DPO significantly lowers the difficulty of scaling up language model training based on human preferences.

\textbf{Limitations \& Future Work.} Our findings open multiple avenues for future investigation. For instance, to what extent does the DPO-trained policy extrapolate to data outside the training distribution in comparison with approaches that learn from explicit reward functions? Preliminary experiments indicate that DPO-trained models show similar generalization abilities to PPO-based alternatives, though extensive evaluation remains necessary. Additionally, it remains an open question whether leveraging self-generated labels from a DPO-trained policy can effectively utilize unannotated prompts. Another pertinent consideration is the manner in which reward over-optimization presents itself within the context of direct preference optimization; in particular, it is unclear whether the minor drop in performance observed in Figure~\ref{fig:dialogue-main}-right serves as an example of this phenomenon. Furthermore, while our experiments extend to models containing up to 6B parameters, scaling DPO to cutting-edge architectures significantly larger in size is a promising research direction. Concerning evaluation protocols, we observed that win rates determined by GPT-4 are influenced by prompt design; investigating optimal strategies for eliciting reliable assessments from automated evaluators remains important. Finally, DPO holds potential utility well beyond language modeling, such as in training generative models operating over alternative data modalities.